\documentclass[a4paper,11pt]{article}
\usepackage[T2A]{fontenc}
\usepackage[utf8]{inputenc}
\usepackage[russian]{babel}
\usepackage{amsmath,amssymb,mathtools}
\usepackage{geometry}
\usepackage{booktabs,tabularx,array}
\usepackage{tikz}
\usepackage{xcolor}
\usepackage{hyperref}
\usepackage{fancyhdr}
\usepackage{microtype}
\usepackage{needspace}
\geometry{left=22mm,right=22mm,top=24mm,bottom=23mm}
\definecolor{accent}{HTML}{008080}
\definecolor{darkaccent}{HTML}{005458}
\definecolor{ink}{HTML}{292E31}
\definecolor{redmark}{HTML}{A3322C}
\hypersetup{colorlinks=true,linkcolor=darkaccent,pdfauthor={Лёвин Артём Александрович},pdftitle={Ревью домашней работы — 13 ноября 2026}}
\setlength{\parindent}{0pt}
\setlength{\parskip}{6pt}
\renewcommand{\arraystretch}{1.38}
\pagestyle{fancy}
\fancyhf{}
\renewcommand{\headrulewidth}{0pt}
\renewcommand{\footrulewidth}{0.25pt}
\fancyfoot[C]{\scriptsize Лёвин Артём Александрович эксклюзивно для Даниила Хрястова}
\fancyfoot[R]{\scriptsize\thepage}
\newcommand{\para}[1]{\par\needspace{3\baselineskip}\textbf{#1}\par\nobreak}
\begin{document}
\thispagestyle{empty}
\begin{center}
\vspace*{30mm}
{\color{accent}\large МАТЕМАТИКА}\par
\vspace{17mm}
{\Huge Ревью домашней работы}\par
\vspace{13mm}
{\color{darkaccent}\Large Метод опорных точек}\par
\vspace{5mm}
{\large Функции и их графики}\par
\vspace{30mm}
{\normalsize 13 ноября 2026 г.}\par
\vspace{15mm}
{\normalsize Автор: Лёвин Артём Александрович}\par
\vfill
{\small Лёвин Артём Александрович эксклюзивно для Даниила Хрястова}\par
\vspace{13mm}
\begin{tikzpicture}
\draw[accent,thin] (0,0) .. controls (1.5,0.8) and (3,-0.7) .. (4.6,0)
  .. controls (6.2,0.8) and (7.7,-0.7) .. (9.2,0);
\end{tikzpicture}
\end{center}
\clearpage
\section*{Сводная таблица}
\vspace{4mm}
\small
\begin{tabularx}{\linewidth}{@{}>{\centering\arraybackslash}p{13mm}p{22mm}X>{\centering\arraybackslash}p{25mm}>{\centering\arraybackslash}p{24mm}@{}}
\toprule
\textbf{№} & \textbf{Тема} & \textbf{Суть ошибки} & \textbf{Ваша запись} & \textbf{Корректная запись}\\
\midrule
\hyperref[task:3]{3} & Экспонента & Несоответствие аргумента функции и числа, подставленного в показатель. & $f(3)=9$ & $f(-3)=9$\\
\bottomrule
\end{tabularx}
\normalsize
\clearpage
\thispagestyle{empty}
\vspace*{\fill}
\begin{center}
{\color{darkaccent}\Huge Разбор ошибок}
\end{center}
\vspace*{\fill}
\clearpage
\phantomsection\label{task:3}
\section*{Задание № 3. Аргумент функции и отрицательная степень}
\para{Текст задания}
По графику функции $f(x)=a^x+b$ восстановить её формулу и найти $f(-3)$. Опорные точки графика: $(0;2)$ и $(-1;3)$. Для показательной функции $a>0$, $a\ne1$.

\para{Ваше решение}
По записи на фотографии видно, что Вы верно получили $b=1$ и $a=\tfrac12$; затем использовали отрицательную степень $-3$ и получили число $9$. Однако рядом с этим вычислением несколько раз написано $f(3)$:
\[
  \underbrace{f(3)}_{\text{записан аргумент }3}
  =\left(\frac12\right)^{-3}+1=9.
\]

\para{Ваш ответ}
В тетради указано $f(3)=9$. Число $9$ является верным значением \emph{запрошенного} $f(-3)$.

\para{Ошибка}
\textcolor{redmark}{\textbf{Ошибка:}} в обозначении функции потерян знак минус у аргумента. Это меняет смысл равенства, даже когда числовое вычисление выполнено правильно.

\para{Где именно возникает ошибка}
\textbf{Последний правильный шаг:} $a=\tfrac12$, $b=1$, следовательно, $f(x)=\left(\tfrac12\right)^x+1$.

\textbf{Первый неверный шаг в записи:} обозначение $f(3)$ рядом с выражением $\left(\tfrac12\right)^{-3}+1$. Подстановка $-3$ соответствует $f(-3)$, а не $f(3)$.

\para{Почему это неверно}
В записи $f(t)$ число $t$ указывает, при каком аргументе вычисляется функция. Если аргумент равен $3$, то по определению надо поставить в формулу именно $x=3$:
\[
  f(3)=\left(\frac12\right)^3+1=\frac18+1=\frac98.
\]
Если аргумент равен $-3$, то необходимо поставить $x=-3$:
\[
  f(-3)=\left(\frac12\right)^{-3}+1=8+1=9.
\]
Значения отличаются. Поэтому нельзя оставлять метку $f(3)$ при вычислении для $x=-3$. Отрицательный показатель имеет особый смысл: при ненулевом основании он означает обратную величину, $a^{-n}=\frac1{a^n}$.

\para{Ключевая идея}
\textcolor{darkaccent}{\textbf{Ключевая идея:}} сначала переписывайте \emph{ровно тот аргумент}, который дан в вопросе, а затем подставляйте это же число в формулу. Знак минус входит в значение аргумента и должен сохраняться.

\para{Исправление от последнего правильного шага}
Оставляем найденную Вами формулу и исправляем обозначение слева:
\[
  \boxed{f(-3)=\left(\frac12\right)^{-3}+1=9.}
\]
Пересчитывать коэффициенты заново не требуется.

\para{Полное правильное решение}
Точка $(0;2)$ лежит на графике, следовательно, $f(0)=2$. Подставляем $x=0$ и $y=2$:
\[
  2=a^0+b=1+b\quad\Longrightarrow\quad b=1.
\]
Точка $(-1;3)$ также лежит на графике. Подставляем её координаты и уже найденное $b$:
\[
  3=a^{-1}+1\quad\Longrightarrow\quad a^{-1}=2.
\]
Так как $a^{-1}=\frac1a$, имеем $\frac1a=2$, откуда $a=\frac12$. Это основание положительно и не равно единице, поэтому допустимо. Формула функции:
\[
  f(x)=\left(\frac12\right)^x+1.
\]
Требуется значение при $x=-3$:
\[
\begin{aligned}
  f(-3)&=\left(\frac12\right)^{-3}+1
  =\frac{1}{(\frac12)^3}+1\\
  &=\frac1{\frac18}+1=8+1=9.
\end{aligned}
\]

\para{Проверка}
Для первой опорной точки $f(0)=1+1=2$; для второй $f(-1)=2+1=3$. Обе точки удовлетворяют восстановленной формуле. Дополнительно $f(3)=\tfrac98$, что подтверждает необходимость различать $f(3)$ и $f(-3)$.

\para{Итог}
\textcolor{darkaccent}{\textbf{Итог:}} правильное значение $\boxed{f(-3)=9}$. Вы нашли нужное число; для математически корректного оформления важно сохранить минус в аргументе на каждом шаге.

\para{Задача на закрепление}
Показательная функция $g(x)=a^x+b$ проходит через точки $(0;2)$ и $(-1;4)$. Восстановите её формулу и найдите $g(-2)$.
\end{document}
